\documentclass[12pt]{article}

% --- Page layout ---
\usepackage{geometry}
\geometry{margin=1in}

% --- Math ---
\usepackage{amsmath,amssymb,amsthm}
\usepackage{mathtools}
\usepackage{enumitem}

% --- Fonts ---
\usepackage{newtxtext}
\usepackage{newtxmath}

% --- References / links ---
\usepackage{hyperref}
\usepackage{cleveref}

% --- Bibliography ---
\usepackage[
    backend=biber,
    style=authoryear,
    sorting=nyt
]{biblatex}
\addbibresource{references.bib}

% --- Theorem environments ---
\theoremstyle{plain}
\newtheorem{theorem}{Theorem}[section]
\newtheorem{lemma}[theorem]{Lemma}
\newtheorem{proposition}[theorem]{Proposition}
\newtheorem{corollary}[theorem]{Corollary}

\theoremstyle{definition}
\newtheorem{definition}[theorem]{Definition}
\newtheorem{example}[theorem]{Example}

\theoremstyle{remark}
\newtheorem{remark}[theorem]{Remark}

% --- Convenience macros ---
\newcommand{\R}{\mathbb{R}}
\newcommand{\N}{\mathbb{N}}
\newcommand{\Z}{\mathbb{Z}}
\newcommand{\C}{\mathbb{C}}

\title{Assorted Notes}
\author{}
\date{\today}

\begin{document}

\maketitle
\tableofcontents

\section{Six Operations}

A good source is \cite{Gallauer2022SixFunctor}.

One important example is the relative point of view as discussed in \cite{Gallauer2022SixFunctor}. One discovers invariants that are easy to compute. An example would be rational points of a variety over a finite field. Then, one finds a corresponding cohomology group associated with this. In this case, it is the $l$-adic cohomology. This idea of a cohomology group was explored by Galois. In \cite{Andre2008Ambiguity}, it is ambiguous what you want to consider given a constant invariant, and this ambiguity is captured by a group.

To find a higher structure of cohomology, one can consider a cohomology group to be something local. To assemble something local to something global, one invents the idea of a $l$-adic constructible derived category. This is a derived category of constructible sheaves. Sheaves, or higher versions of sheaves with better combinatorics, give the right properties.

To invent the notion of global to local correspondences, one needs some form of restriction on coverings. Hence, one invents the notion of a presheaf as a contravariant functor of space, most a combinatorial model such a simplicial set. Then, this is valued in an $(\infty,1)$-groupoid to capture automorphisms. \cite{nLabInfinityOneSheaf} gives a definition.

The more basic example is a presheaf as a contravariant functor from the category of open sets of a topological space to the category of sets. One loses the ability to capture automorphisms. Set valued presheaves cannot capture automorphisms. For sheafification, one wants separability or the representatives being unique on overlaps, and descent, which is a condition for covering sieves. For descent, or gluing, Cech cocyles are effective with an equaliser diagram. The purpose of descent is to capture the idea of a local to global correspondence.

We now work in the category of morphisms of schemes. Why should I care about adjoint functors? This morphism induces adjoint functors which are ordinary pullback and ordinary pushforward between categories of sheaves.
\section{Algebraic geometry}

One invents the notion of a diagonal map to understand separability. The diagonal map is the prototypical example of self intersection in a space.

Separability means diagonal pullback is a closed immersion.
\input{notes/stacks}
\input{notes/category}
\input{notes/chromatic-homotopy}
\input{notes/equivariant}
\section{K-theory}

I want to talk about using trace maps, then TCC and THH to assemble algebraic K-theory.
\section{Mathematical physics}

Good sources for this is D. B. Skinner's Advanced Quantum Field Theory course notes.
\section{Prismatic F-gauges}

Jacob Lurie's introductory talk is most accessible for this.

\printbibliography

\end{document}